% =====================================================================
\section{Metodología de ingeniería de requisitos}
\label{sec:metodologia}

\subsection{Marco normativo aplicado}

El ERS se estructura conforme a ISO/IEC/IEEE 29148:2018, que define el contenido de una
especificación de requisitos de software y las características exigibles a un requisito individual y
al conjunto \cite{iso29148}. La calidad del producto documental se interpreta con el modelo de
ISO/IEC 25010:2023, del que la \S\ref{sec:metricas} deriva seis métricas operativas
\cite{iso25010}. El encuadre de procesos de ciclo de vida proviene de ISO/IEC/IEEE 15288:2023 para
sistemas y de ISO/IEC/IEEE 12207:2017 para software \cite{iso15288, iso12207}, que sitúan la
definición de requisitos de las partes interesadas y el análisis de requisitos del sistema dentro de
los procesos técnicos.

Los fundamentos conceptuales siguen a Pohl \cite{pohl2025} y al SWEBOK v4.0 \cite{swebok4}, y la
lista de verificación de calidad de la especificación se apoya en el temario CPRE Foundation Level
del IREB \cite{ireb31}. La gestión del proyecto integrador se organiza con los principios del PMBOK
7.ª edición \cite{pmbok7}.

Para los componentes de inteligencia artificial concurren dos marcos adicionales: el Reglamento
(UE) 2024/1689, que clasifica los sistemas por nivel de riesgo e impone obligaciones proporcionales
\cite{aiact2024}, y la Ley Orgánica de Protección de Datos Personales del Ecuador, que rige el
tratamiento de los datos personales que el sistema recoge \cite{lopdp2021}.

\subsection{Fundamento teórico}

\subsubsection*{Validación frente a verificación}

La distinción es anterior a cualquier técnica y condiciona qué se puede concluir de una inspección.
Verificar es comprobar que el documento está bien construido conforme a un estándar; validar es
comprobar que especifica lo que las partes interesadas necesitan \cite{pohl2025}. Un ERS puede
superar la verificación y fallar la validación: sus requisitos serían completos, consistentes y
verificables, y describirían un sistema que nadie pidió.

ISO/IEC/IEEE 29148:2018 sitúa la validación en su cláusula 6.4 y enumera las características
exigibles a un requisito individual ---necesario, apropiado, no ambiguo, completo, singular,
factible, verificable, correcto y conforme--- y al conjunto ---completo, consistente, factible,
comprensible, capaz de ser validado \cite{iso29148}. Las seis métricas de la
\S\ref{sec:metricas} operacionalizan seis de esas características; las restantes no se midieron
porque no admiten un conteo objetivo sin introducir juicio del auditor.

\subsubsection*{Técnicas de validación y el método Fagan}

Las técnicas de validación se ordenan por grado de formalidad: revisión informal, recorrido
(\emph{walkthrough}), inspección formal, prototipado y listas de verificación \cite{pohl2025,
swebok4}. El método de Fagan es la inspección formal canónica y aporta tres elementos que las
técnicas menos formales no tienen: roles diferenciados con responsabilidad distinta, preparación
individual obligatoria antes de la reunión, y prohibición expresa de discutir soluciones durante la
sesión \cite{fagan1976}.

Los tres elementos responden al mismo problema. Sin roles, todos leen igual y detectan lo mismo. Sin
preparación individual, el primer comentario ancla al grupo. Y sin la prohibición de discutir
soluciones, la sesión deriva hacia el diseño y deja de recorrer el documento. Las métricas de
inspección ---densidad de defectos, distribución por tipo y severidad, tasa de detección por
inspector y esfuerzo--- permiten evaluar el proceso y no solo su resultado.

\subsubsection*{Gestión de requisitos: línea base, trazabilidad y cambio}

La gestión de la configuración de un ERS descansa en la noción de línea base: una versión aprobada y
congelada que sirve de referencia y a partir de la cual todo cambio se tramita \cite{swebok4}. Sin
línea base no hay control de cambios posible, porque no existe el estado respecto del cual algo
cambia.

La trazabilidad se distingue en \emph{pre-traceability} ---del requisito hacia su origen--- y
\emph{post-traceability} ---del requisito hacia los artefactos que lo realizan--- \cite{gotel1994}.
Ambas direcciones responden preguntas distintas: la primera, por qué existe un requisito; la segunda,
qué hay que tocar si cambia. La matriz de la \S\ref{sec:gestion} las implementa como columnas hacia
atrás y hacia adelante de una misma cadena.

El proceso de cambio ---solicitud formal, análisis de impacto, decisión de un órgano colegiado,
implementación y cierre--- sigue el modelo de control integrado de cambios del PMBOK
\cite{pmbok7}. La literatura de trazabilidad advierte de que el análisis de impacto estimado sin
recorrer la matriz es la causa más frecuente de subestimación del costo de un cambio
\cite{clelandhuang2012}.

\subsubsection*{Herramientas CASE y trazabilidad en entornos ágiles}

Las herramientas CASE de gestión de requisitos se clasifican por su posición en el ciclo de vida
---\emph{upper}, \emph{lower}, integradas--- y por el conjunto de funcionalidades esperables:
repositorio, atributos tipados, trazabilidad multinivel, \emph{baselining}, control de cambios,
colaboración e integración con el control de versiones \cite{wiegers2013}.

En entornos ágiles la trazabilidad no desaparece: se reconstruye por convención sobre el mensaje de
\emph{commit} y la vinculación de ramas a la incidencia \cite{clelandhuang2012}. El formato de
historia de usuario \cite{cohn2004} y el criterio de aceptación en notación Gherkin \cite{smart2014}
cumplen el papel que en la especificación documental cumplen la ficha de requisito y su criterio de
verificación. El proyecto mantiene ambas representaciones, y la \S\ref{sec:gestion} cuantifica el
costo de esa convivencia.

\subsubsection*{Ingeniería de requisitos para sistemas con aprendizaje automático}

Especificar un componente cuyo comportamiento se induce de datos plantea problemas que la ingeniería
de requisitos clásica no aborda. El resultado es probabilístico, de modo que un criterio de
aceptación binario no aplica y hay que sustituirlo por una métrica con umbral y conjunto de
evaluación declarado. La calidad depende del conjunto de datos tanto como del modelo, lo que convierte
las propiedades de los datos en requisitos. Y aparecen categorías de calidad nuevas ---explicabilidad
y ausencia de discriminación--- que ningún modelo de calidad tradicional contemplaba
\cite{vogelsang2024, chazette2021}.

A esos requisitos técnicos se suman los normativos. El Reglamento (UE) 2024/1689 adopta un enfoque
basado en riesgo: prohíbe determinadas prácticas, impone obligaciones reforzadas a los sistemas de
alto riesgo enumerados en su Anexo III, y establece obligaciones de transparencia cuando el sistema
interactúa directamente con personas \cite{aiact2024}. La normativa ecuatoriana de protección de
datos personales rige el tratamiento de los datos que alimentan el modelo y los derechos del titular
sobre ellos \cite{lopdp2021}. La \S\ref{sec:ia} aplica ambos marcos al SIMPA y declara la
clasificación resultante con su razonamiento.

\subsection{Proceso seguido a lo largo de las cinco unidades}

\begin{table}[H]
\centering\footnotesize
\caption{Proceso de ingeniería de requisitos por unidad}
\label{tab:proceso}
\begin{tabular}{|C{1.4cm}|L{4.0cm}|L{4.6cm}|L{4.6cm}|}
\hline
\rowcolor{grisclaro}\textbf{Unidad} & \textbf{Actividad} & \textbf{Técnicas empleadas} & \textbf{Artefacto producido}\\ \hline
I & Fundamentos y planificación & Análisis del dominio, identificación de partes interesadas, matriz poder/interés & Plan de elicitación y mapa de partes interesadas\\ \hline
II & Elicitación y análisis & Entrevista semiestructurada, cuestionario, observación de campo, análisis de documentos & Trece evidencias (\id{EV-01} a \id{EV-13}) y catálogo de requisitos brutos\\ \hline
III & Especificación y modelado & Fichas de requisito conforme a 29148, casos de uso, modelado i*, UML, diagramas de flujo de datos & ERS de las Entregas 1B y 2A; modelos UML e i*\\ \hline
IV & Validación y gestión & Inspección formal tipo Fagan, Change Control Board simulado, trazabilidad en Jira, línea base con Git & Registro de 33 defectos, tres RFC, acta del CCB, línea base \id{baseline-v1.1}\\ \hline
V & Integración y auditoría & Auditoría con métricas de 25010, re-inspección, especificación de componentes de IA & Instrumento de auditoría, matriz extremo a extremo, \S9 del ERS\\ \hline
\end{tabular}
\end{table}

\subsection{Decisiones metodológicas y su justificación}

Cuatro decisiones condicionaron el resultado y merecen justificarse, porque en cada una había una
alternativa razonable que se descartó.

\subsubsection*{Elicitación con instrumentos mixtos y no solo con entrevistas}

La elicitación combinó entrevista semiestructurada con la administración, cuestionario al personal de
campo y observación directa. La razón es una asimetría documentada en \id{EV-07}: la administración
declara con soltura los procesos formales, pero el personal de campo ---sesenta y dos personas con
competencia digital declarada como baja--- no responde bien al formato de entrevista. Basar la
elicitación solo en entrevistas habría producido un catálogo de requisitos que refleja cómo la
administración cree que funciona el campo, y no cómo funciona.

El cuestionario ampliado \id{EV-12}, con $n=62$, es el que aportó los dos datos que más restricciones
generaron: que el 11,3\,\% del personal no dispone de teléfono inteligente y que el 74,2\,\% no tiene
señal permanente en el lote. Ambos datos contradijeron suposiciones que el equipo había dado por
válidas y obligaron a especificar \id{RF-35} (registro delegado) y \id{RD-02} (operación sin
conexión).

\subsubsection*{Priorización con tres técnicas y no con una}

La priorización combina MoSCoW, el modelo de Kano y el cálculo WSJF. MoSCoW por sí solo produce
categorías sin orden interno: veinticuatro requisitos \emph{Must} no dicen cuál construir primero.
Kano aporta la distinción entre lo que satisface y lo que se da por supuesto. WSJF aporta el orden
dentro de la categoría, ponderando valor de negocio, urgencia, reducción de riesgo y tamaño.

El resultado de combinarlas no fue redundante. \id{RF-35} ---el registro delegado para quien no tiene
teléfono--- encabeza el orden de implementación con un WSJF de 7,00 pese a no ser el requisito de
mayor valor de negocio: lo que lo eleva es la reducción de riesgo, porque sin él el 11,3\,\% del
personal quedaría fuera del sistema y el registro de avance sería incompleto por diseño.

\subsubsection*{Inspección formal y no revisión informal}

La Unidad IV empleó el método Fagan con roles diferenciados, preparación individual previa registrada
con marca temporal y prohibición expresa de discutir soluciones durante la sesión \cite{fagan1976}.
La alternativa ---una revisión conjunta del documento--- es más barata y detecta menos: sin
preparación individual, la sesión se convierte en una lectura colectiva donde el primer comentario
ancla al resto del grupo.

El resultado empírico sostiene la decisión: treinta y tres defectos únicos sobre sesenta y tres
páginas efectivas, con una densidad de 0,52 defectos por página, y ningún rol por debajo del mínimo
de seis defectos preparados.

\subsubsection*{Formato BDD para los criterios de aceptación}

La Unidad V reescribió treinta y siete criterios de verificación en formato Dado--Cuando--Entonces
\cite{smart2014}. La alternativa era conservarlos en prosa, que es como estaban y que no es
incorrecto. La razón del cambio es que un criterio en prosa admite que dos evaluadores lleguen a
veredictos distintos sobre el mismo resultado, mientras que la plantilla obliga a fijar contexto,
evento y resultado observable. La verificabilidad, tal como la define ISO/IEC/IEEE 29148:2018, no es
que exista un criterio sino que el criterio permita decidir \cite{iso29148}.

Conviene precisar el alcance de ese cambio, porque es fácil sobreinterpretarlo: no se especificó
ningún requisito nuevo ni se modificó ningún umbral. Cambió la forma de enunciar la condición, no su
contenido.

